\section{Discussion and limitations}\label{sec:discussion}

\paragraph{What the sums-only restriction costs}
\clamp{} gives up two things to stay compatible with secure aggregation. First, its robustness is a norm bound, not a filter: an adversary that stays within the bound still shifts the average by up to $2fC/(K-f)$ per round (Theorem~\ref{thm:update}), and rules that inspect individual updates can do better against attacks whose updates look unlike honest ones. Second, clipping at a fixed $C$ slows honest training when honest updates are larger than $C$, and may contribute to the small utility gap between ClipMean and FedAvg-GAT on some datasets (Table~\ref{tab:e1_main_compact}); a looser clip did not improve AUC in our ablation, so we have not isolated this cost. An adaptive, privately estimated bound is a natural extension, at the price of another channel to account for and protect.

\paragraph{Small graphs}
The fairness controller is only as good as its gap estimate, and on small graphs the validation estimate is a poor proxy for the test gap. On German, the 250 validation nodes are split over five clients, so a client typically holds 15 to 35 validation nodes per group. Statistics noise adds to this, but it is not the main problem: switching the noise off does not make control effective on NBA or Facebook. Larger or pooled validation sets, or a fairness criterion with lower variance, would be needed. Theorem~\ref{thm:stats} makes the same point quantitatively: its bound is weak when honest validation mass is small relative to $f\ncap$.

\paragraph{What is and is not private}
Theorem~\ref{thm:dp} is an attribute-level guarantee for the statistics channel only. The trained model, and therefore the aggregate update, still depends on the training nodes' attributes through the fairness penalty, and the probe in Section~\ref{sec:privacy} shows how much of that dependence is measurable. Making the whole pipeline differentially private at the client level would need noise on the updates, which we evaluate only as an ablation.

\paragraph{Simulation of the cryptographic layer}
We simulate verifiable secure aggregation in plaintext. The guarantees assume a deployment in which norm and range proofs are enforced cryptographically and enough honest clients participate for secure aggregation to hide individual inputs. We do not measure the cost of that layer; systems such as RoFL~\citep{lycklama2023rofl} report it.

\paragraph{Scope of the fairness notion}
We control one binary sensitive attribute and the demographic-parity gap. Equal opportunity needs label-conditioned statistics, which fit the same template (four more sums per release; the $\ell_2$ sensitivity grows from $2$ to $2\sqrt{2}$ under the attribute-neighbouring relation), but we have not evaluated it. Multi-valued attributes need one multiplier per constraint.
